\section{Wide focus (Theorem~\ref{thm:wide})}
\label{sec:wide}

\subsection{Statement}

\begin{proposition}[wide focus, line form] \caproofst{App.~\ref{app:C}}
\label{prop:wide}
Let $\ell\ge960.8$.
Let a state satisfy $\mu>0$, $\sigma^3\le r\le\tfrac14$, $8\ell+100\le B<\ell^4$, $b\le\tfrac12$, $\kappa>0$, and let $P>0$ satisfy (H$'$).
Then for every $a\in(0,4r/\sigma]$ the option chosen by the rule of Sec.~\ref{sec:narrow} (same constants $c_1$, $c_T$, $c_{\rm ch}$, $a_s$) satisfies $\Lambda_k(a)\ge0$ with $\gamma=\Gamma_\ell$.
\end{proposition}

Theorem~\ref{thm:wide} follows exactly as Theorem~\ref{thm:narrow} does (Lemma~\ref{lem:lines} for $G>0$, Lemma~\ref{lem:G0} for $G=0$).
For $\ell\ge960.8$, $\ell^4\ge8.5\times10^{11}$, so the proposition overlaps Proposition~\ref{prop:narrow} and covers the band $2.5\times10^{10}\le B<50(8\ell+100)$ that Proposition~\ref{prop:narrow} leaves for $\ell>6.25\times10^7$.
The upper end $\ell^4$ is used in one step only, (b.5).

The focus is wide: $\tau=x/b\ge5\times10^{11}/B$, and the focus spans many periods.
What makes the proof of Sec.~\ref{sec:narrow} work here is that it never uses the smallness of $b$ directly.
The lower end $B\ge2.5\times10^{10}$ entered only through constants, and the lower end $B\ge50(8\ell+100)$ only through the uniform bounds $q\le0.0026$, $z_2\le0.01$.

\subsection{The lines down to $B=8\ell+100$}

For the lines to be lower bounds (Lemma~\ref{lem:lines}) we need $z_0^2\le1$, $1-q>0$ and one constant of the main-term identity, which was established for $y\le10^{-6}$ and is needed here for $y=(1-b)^2/(2B)\le1/(2(8\ell+100))<6.5\times10^{-5}$.

\begin{lemma}[the constant $M_4$] \status{(i)--(ii) proof; (iii) proof, computer-assisted}
\label{lem:M4}
Let $f(u,y)=e^{-u/2}\bigl[\tfrac{1+u}2(1+u\cosh\sqrt{yu})+\cosh\sqrt{yu}\bigr]$ and $M_4(y)=\sup_{u\ge0}f(u,y)$.
(i) The quantity $\Xi$ of the main-term identity satisfies $|\Xi|\le\Qb M_4(y)\,y$.
(ii) $M_4$ is non-decreasing in $y$.
(iii) $M_4(10^{-4})\le2.0475\le2.07$.
\end{lemma}

\begin{proof}
Put $c=\sqrt{yu}$.
(i) $1-e^{-y/2}\cosh c\le1-e^{-y/2}\le y/2$, and $e^{-y/2}\cosh c-1\le\cosh c-1\le(c^2/2)\cosh c$; so $|1-e^{-y/2}\cosh c|\le(y/2)(1+u\cosh c)$. With $ye^{-y/2}\cosh c\le y\cosh c$ the single term is at most $y f(u,y)$, and the sum over harmonics with $\sum q_n\le\Qb$ gives (i).
(ii) For fixed $u$, $\cosh\sqrt{yu}$ is non-decreasing in $y$ and enters with positive coefficients.
(iii) On $u\in[0,60]$, pieces of width $0.005$ ($[0,20]$) and $0.05$ ($[20,60]$) bound $f$ by its decreasing and increasing parts at the ends of each piece, in interval arithmetic; the largest bound is $2.04742$. For $u\ge60$, $f\le e^{-0.4987u}((1+u)^2/2+1)$, which decreases and is below $1.9\times10^{-10}$ at $u=60$.
\end{proof}

All other places where $y$ enters (the background lines, the pair lemma, the lines B1w) hold for every $y\le1$ \proofst.
So the lines are lower bounds for $B\ge8\ell+100$ with the constant $2.07$ unchanged.

\subsection{Proof: the same steps, other constants}

The steps (F), (R), Case~I, (c), (II-ch), (II-s), (m1)--(m4) of Sec.~\ref{sec:narrow} are used as they stand.
On a piece $[\ell_a,\ell_b]$ of $\ell$, with $B_m=8\ell_a+100$ the lower bound of $B$ on the piece, five constants are replaced:
\begin{itemize}
\item $q\le(1+\eta)/(8(1-\eta))$ and $z_2\le\tfrac12$ (instead of $0.0026$ and $0.01$);
\item $\mathrm{es}=0.25/B_m$ in (R1) and (R3); at $\ell=960.8$ the sum of (R1) is $0.50042\,\mathrm{es}$;
\item $e_0=2/\sqrt{B_m}+\mathrm{es}$ in (R4) ($e_0=0.022697$ at $\ell=960.8$, against $1.26\times10^{-5}$ in Sec.~\ref{sec:narrow});
\item $B_m$ inside $\omega$, $c_S$, $c_{1F}$ and $\bar\varrho$;
\item in (m4), $B<\ell^4$ closes the argument by (b.5) alone, so the conditions (N), (N0) and (F) are not used.
\end{itemize}
Some numbers of Sec.~\ref{sec:narrow} are written into the formulas and hold only for $z_2\le0.01$ or $q\le0.0026$.
They were found by reading the source line by line and replaced by arguments: $k_f\ge0.999$ in (b.5) becomes $k_f\ge(1-b)^2(1-4b/(1-b)-z_2/12)$ (with $\ell^4\sigma^2\le10^{-100}$ checked), $1/\sqrt{c_0}\ge0.999$ becomes $(m_2(1-q))^{1/4}$, the factor $2.1$ in $\mathrm{A2k}\le2.1b^2\mathrm{cKz}$ is checked as a condition (its value is $2.0915$), and three smallness claims of (II-ch), (II-s) and Case~I are checked as conditions.
Two independent counts of the written-in numbers find none missing.
At $z_2\le\tfrac12$ some constants move: in Case~I with Tz-B2, $k_f\ge0.4162$ and $\mathrm{A2k}\le0.0482\,\mathrm{cKz}$ at $\ell=960.8$.

\subsection{The computer-assisted part}

The constants and conditions are evaluated by the same interval arithmetic as in Sec.~\ref{sec:ca}, built from the source of the computation of Sec.~\ref{sec:narrow} with only the quantities above turned into arguments (with the old values it reproduces that computation to $2.2\times10^{-36}$), on $[960.8,10^6]$ cut into $6{,}952$ pieces and on the tail \caproofst{App.~\ref{app:C}}.
$35$ conditions are checked on each piece; none fails on any of the $6{,}953$ pieces.
The smallest margins are at $\ell=960.8$: Case~I with Tz-B2 $0.0023301$, Case~I with Fz $0.07874$, (II-ch) $0.038309$, and the absorption of the constant lines uses $0.96287$ of $e_0$.
Each of four deliberately wrong versions ($e_0:=\mathrm{es}$; $z_2>1$; the rule constant $c_1=0.95$; $B_m=100$) fails on all four test pieces.

\emph{Where the end $8\ell+100$ comes from.}
The constants alone go much lower: all conditions hold with $B_m=10^3$, and Case~I with Fz fails first at $B_m=300$.
The end $8\ell+100$ is the end of the lines: below $B=(8\ell+100)/2$ the ring radius $K^2=8\ell+100$ violates $z_0^2=K^2y\le1$.
The searches of Sec.~\ref{sec:break} see this: below that end the sum of the lines becomes negative for every $\gamma$, which is not a counterexample, because there the lines are not lower bounds.
This is where Theorem~\ref{thm:verywide} takes over.

\emph{Picture.}
A wide focus spans many periods, so the brightness of the ripple is spread over it ($\kappa\approx1$), and the phase turns appreciably between the centre and the edge of the focus.
The chirp gain $\chi G$ grows like $\kappa\sigma/b$, that is, like the width of the focus, and exceeds the apparent darkening of the ripple by orders of magnitude.
So the hardest states of this region are its narrowest ones, at the border with Theorem~\ref{thm:narrow}; widening the focus only helps (Sec.~\ref{sec:break}).
